\documentclass[11pt]{article}

\usepackage[margin=25mm]{geometry}

\usepackage{amsmath,amssymb,fontspec,unicode-math}

\IfFontExistsTF{DejaVu Serif}{\setmainfont{DejaVu Serif}}{\setmainfont{Latin Modern Roman}}

\IfFontExistsTF{Latin Modern Math}{\setmathfont{Latin Modern Math}}{\setmathfont{DejaVu Math TeX Gyre}}

\usepackage{xurl,hyperref}

\setlength{\parindent}{0pt}\setlength{\parskip}{6pt}

\begin{document}

\title{DiffLiG: Diffusion-enhanced Liquid Graph with Attention Propagation for Grid-to-Station Precipitation Correction}

\author{ELVIA preliminary AI analysis}\date{21 October 2026}\maketitle

Preliminary assessment: Minor.\par

Independent scientific review remains pending. Author adoption has not been requested.

The official implementation uses incompatible geographic coordinates, retains closed edges after softmax, has a constant consensus score, discards grid history and exceeds the reported parameter count in one submodule alone. Published forecast effects remain unmeasured.

\section*{Review scope}

Claim-focused reading of the original paper and relevant author-released passages; direct mathematical or static-source reasoning. No author program, earlier verifier, model or benchmark was executed during this check.

\section*{Problem 1: Geographic degrees and normalized coordinates are conflated}

Source: Section 3.4 and Algorithm 1, PDF pp.5–6; Appendix J, PDF pp.33–34; MixData.py and Haversine layer methods

Assessment: supported; impact: minor

\subsection*{Source claim (paraphrase)}

The paper defines physical Haversine distances and geographic nearest neighbors; the release feeds min–max-normalized coordinates into a degree-based Haversine formula.

\subsection*{Analysis}

Points \((100^\circ,30^\circ)\) and \((100.125^\circ,30^\circ)\) are about \(12.0372\) km apart with Earth radius \(6371\) km. After longitude bounds \(70\)–\(140\) and latitude bounds \(15\)–\(55\), the same formula on normalized values gives about \(0.198558\) km, reversing the predicate at threshold \(0.5\). Coordinate-Euclidean ranking also differs: from \((100^\circ,50^\circ)\), candidates \((101.4^\circ,50^\circ)\) and \((100^\circ,51^\circ)\) have coordinate distances \(1.4,1\), but Haversine distances approximately \(100.063,111.195\) km, reversing the nearest neighbor.

\subsection*{Finding and scope}

The physical graph operations differ from the description. Coordinate normalization for learned embeddings remains legitimate; benchmark consequences are unresolved.

\subsection*{Proposed correction}

Keep raw coordinates for physical distances and normalized coordinates for embeddings; rank the intended neighbors using the declared physical metric.

A proposed correction is retained where stated; independent verification and author adoption remain unestablished.

\section*{Problem 2: Multiplying a logit by zero does not remove its edge}

Source: Section 3.4, PDF p.5; GNN\_Layer\_External.message, lines 87–99; internal counterpart

Assessment: supported; impact: minor

\subsection*{Source claim (paraphrase)}

The paper multiplies normalized attention by gate and decay; the release applies softmax after multiplying the raw logits.

\subsection*{Analysis}

For logits \((-1,8)\), decays \((0.5,0.2)\), gates \((1,0)\), the source normalizes logits \((-0.5,0)\), yielding weights approximately \((0.377541,0.622459)\). The closed edge receives positive, larger mass, while the paper’s factorized expression has second coefficient exactly \(0\). All-closed groups become uniform. A hard threshold comparison converted to floating point also supplies no ordinary gradient to the threshold, despite its parameter declaration.

\subsection*{Finding and scope}

Closed edges survive aggregation and the hard threshold does not implement the described learned gate path.

\subsection*{Proposed correction}

Apply the declared gate/decay after normalized attention with explicit neighborhood and empty-group conventions. A masked-softmax alternative requires its own specified normalization and gate-gradient rule.

A proposed correction is retained where stated; independent verification and author adoption remain unestablished.

\section*{Problem 3: The consensus score is constant and aggregation crosses examples}

Source: Section 3.6, PDF pp.6–7; selector.py, forward and score reductions

Assessment: supported; impact: minor

\subsection*{Source claim (paraphrase)}

The selector is described as attention across candidates, assigning a candidate-specific consensus score.

\subsection*{Analysis}

Its input \((N,M,d)\) is permuted to \((M,N,d)\) for default sequence-first attention, so stations are the sequence and candidates are batches. Averaging both station dimensions of a row-stochastic attention matrix gives \(M^{-2}\sum_{u,v}A_{nuv}=1/M\) for every candidate in evaluation mode. Thus candidates \((0,1,3)\) and \((4,2,1)\) both score \(1/3\). Other selector branches remain active. Their unsegmented averaging can also make one example depend on batch mates: candidate scores \((0,2)\) give softmax \((0.119203,0.880797)\), but adding another example with \((10,-10)\) makes the pooled logits \((5,-4)\), reversing the first example’s mixture.

\subsection*{Finding and scope}

The claimed consensus branch cannot distinguish candidates at evaluation, and global pooling can couple separate forecast examples.

\subsection*{Proposed correction}

Use the intended candidate/station attention layout and a nonconstant, explicitly defined consensus reduction; segment independent example-level score aggregation.

A proposed correction is retained where stated; independent verification and author adoption remain unestablished.

\section*{Problem 4: The grid branch discards every earlier timestep}

Source: Figure 2, PDF p.4; GNN\_Layer\_External.forward, line 66

Assessment: supported; impact: minor

\subsection*{Source claim (paraphrase)}

The described grid encoder attends over temporal input, while its active forward path selects only the final column.

\subsection*{Analysis}

Two grid nodes with histories \((1,2,3,4)\), \((5,6,7,8)\) and alternatives \((1001,1002,1003,4)\), \((1005,1006,1007,8)\) both map to \((4,8)\). Holding other inputs and randomness fixed, all downstream values in this branch agree; derivatives with respect to discarded times are zero. A temporal-attention module is constructed but unused in the active path. The separate station-history encoder remains active.

\subsection*{Finding and scope}

The release cannot use the claimed grid temporal history. This does not mean the entire model has no temporal information.

\subsection*{Proposed correction}

Either implement/document the intended temporal encoder or state that only the final gridded field is used. No replacement architecture is established as restoring results.

A proposed correction is retained where stated; independent verification and author adoption remain unestablished.

\section*{Problem 5: One configured submodule exceeds the reported full-model count}

Source: Table 2, PDF p.7; DiffLiG\_Launcher.py, line 31; SimpleGraphDiffusion constructor

Assessment: supported; impact: minor

\subsection*{Source claim (paraphrase)}

Table 2 reports \(49{,}361\) full-model parameters, while the launcher uses hidden width \(h=128\).

\subsection*{Analysis}

The active diffusion constructor contains three \(h\to h\) graph convolutions, two \(2h\to h\) convolutions, a \(2\to h\) position projection, an \(h\to h\) output projection and an \(h\to1\) auxiliary head. Their standard total is \(8h^2+10h+1=132{,}353\). Even the five graph-convolution matrix weights alone contain \(7h^2=114{,}688\), independent of bias or auxiliary-head conventions. Ensemble reuse does not multiply this count.

\subsection*{Finding and scope}

The released configuration cannot be the reported parameter-count configuration. The exact table-producing configuration is unknown.

\subsection*{Proposed correction}

Identify the actual experiment configuration/checkpoint and count its instantiated parameters; no guessed replacement full-model count is supplied.

A proposed correction is retained where stated; independent verification and author adoption remain unestablished.

\section*{Limits}

Preliminary AI assessment; independent human scientific adjudication remains pending.\par

Analytic witnesses and static release checks do not remeasure the published benchmark results.\par

\section*{Source identity}

NeurIPS 2025 official proceedings PDF (retained local copy)

\url{https://proceedings.neurips.cc/paper_files/paper/2025/file/db7c4a8a0c1bf8b405f21c783bbf9eae-Paper-Conference.pdf}

Reviewed PDF SHA-256: \nolinkurl{347027c7345982c94a304d847f15c11225ac95c8dce9fab33a4a1d82fcfd1416}

The SHA-256 identifies the exact retained PDF inspected in this rapid source/argument check. Original supplied report bytes are retained locally. The original PDF retrieval timestamp is taken from matching retained metadata.

\end{document}
